\documentclass{article}
\usepackage{graphicx} % Required for inserting images
\usepackage{hyperref}
\usepackage[margin=1in]{geometry}

\newcommand{\vs}[1][1]{         %THere is one argument, and the default is 1.
    \vspace{#1em}
}

\title{Tensor Networks, Quantum Computation and Information Notes}

\author{Kaustubh Singh}

\begin{document}
% Make the cliackable references a different colour TODO
% If the reference is open source, it could be saved in a repo, or a local pointer to it could be saved as a references.
% The arxiv number can be used as a pointer to the locally stored file. you can refer to those instead of the href. 

\maketitle

\begin{center}
    These notes will contain my attempts at exploring new topics in quantum computation and information. These will contain references to books, webpages, papers and articles which I found interesting, along with possibly hyperlinks to my own implementations in certain cases. 
\end{center}

\section{Quantum information}

\href{https://felixleditzky.info/selected_references.html#intros}{Selected references to topics in mathematics and quantum information theory} by \href{https://felixleditzky.info/}{Felix Leditzky}. A webpage with a very wide variety of references. 

\section{Tensor networks}
% TODO Include the papers in the bibliography as references.
\href{https://arxiv.org/abs/2205.10345}{Tensor Network Algorithms: a Route Map
} by Mari Carmen Bañuls.

% Get rid of the small indent before the new line starts after a \vs 

\vs

\href{https://arxiv.org/abs/cond-mat/0701428}{From density-matrix renormalization group to matrix product states} by Ian P. McCulloch.

\vs

\href{https://arxiv.org/abs/1910.00257}{Review on novel methods for lattice gauge theories
} by Mari Carmen Bañuls, Krzysztof Cichy.

\vs 

\href{https://arxiv.org/abs/quant-ph/0511070}{Classical simulation of quantum many-body systems with a tree tensor network
} by Yaoyun Shi, Luming Duan, Guifre Vidal. This is TEBD with TTNs. 

\vs 

\href{https://arxiv.org/abs/2609.16520}{Truncated hybrid tensor networks for distributed quantum simulation} by Yong Liu, Guangyao Huang, Weixu Shi et. al. This is the hybrid TEBD algorithm.

\section{Collective neutrino oscillations}
\href{https://arxiv.org/abs/2301.00342}{Many-body collective neutrino oscillations: recent developments
} by Amol V. Patwardhan, Michael J. Cervia, Ermal Rrapaj, Pooja Siwach et. al.

\vs

\href{https://arxiv.org/abs/2305.01150}{Quantum information and quantum simulation of neutrino physics} by A. B. Balantekin, Michael J. Cervia, Amol V. Patwardhan, Ermal Rrapaj, Pooja Siwach.
\vs 

\href{https://arxiv.org/abs/2202.01865}{Collective neutrino oscillations with tensor networks using a time-dependent variational principle} by Michael J. Cervia, Pooja Siwach et. al.

\section{Quantum communication}

\href{https://arxiv.org/abs/2506.20484}{   An adversary bound for quantum signal processing
} by Lorenzo Laneve.

\end{document}
